%%═══════════════════════════════════════════════════════════════════
%% Lecture 02 — Invariance of the Spacetime Interval
%% 77401 — Analytical Electrodynamics
%% Date:
%% Lecturer: Prof. Michael Smolkin
%%═══════════════════════════════════════════════════════════════════

\section{Lecture 02 — Invariance of the Spacetime Interval and the
  Lorentz Group}\label{sec:lec02}

\begin{center}
\begin{tabular}{ll}
  \textbf{Date:}\quad 24/03/2026   & \\
\end{tabular}
\end{center}
\hrule\vspace{1.5em}
\subsection*{Overview}

This lecture establishes the central kinematic fact of special
relativity: the \emph{spacetime interval} $ds^2$ is invariant under
passage between inertial frames.  Starting from the specific result
of Lecture~01 — that a signal propagating at the maximum speed $c$
has $ds^2 = 0$ in every frame — Smolkin argues, using homogeneity,
isotropy, and a three-frame consistency condition, that $ds^2$ is
proportional to $ds'^2$ for \emph{any} pair of events, and that the
proportionality constant must equal unity.  The lecture concludes by
connecting the interval to the Lorentz group: the Lorentz
transformations are precisely the linear maps that preserve $ds^2$,
in exact analogy with orthogonal transformations preserving the
Euclidean distance $d\ell^2 = dx^2 + dy^2 + dz^2$.  The next
lecture will derive the explicit form of the Lorentz transformation
from this invariance requirement.

%%──────────────────────────────────────────────────────────────────
\subsection*{Definitions}

\dfn{Spacetime event}{An \emph{event} is a point in spacetime specified by a coordinate
quadruplet $(t, \mathbf{r})$ or, in component form, $(t, x, y, z)$.}

\dfn{Spacetime interval}{Given two events $(t_1, \mathbf{r}_1)$ and $(t_2, \mathbf{r}_2)$ in
an inertial frame $S$, the \emph{spacetime interval} (squared) is
\begin{equation}
  \Delta s^2 = c^2(t_2-t_1)^2 - |\mathbf{r}_2-\mathbf{r}_1|^2.
  \label{eq:lec02:interval}
\end{equation}
In the infinitesimal form:
\begin{equation}
  ds^2 = c^2\,dt^2 - d\mathbf{r}^2.
  \label{eq:lec02:interval_inf}
\end{equation}}

\dfn{Lorentz transformations}{The \emph{Lorentz transformations} are the set of linear maps between
inertial frames that leave the spacetime interval
\eqref{eq:lec02:interval} invariant.  They form a group — the
\emph{Lorentz group} $O(3,1)$ — under composition, in complete
analogy with the orthogonal group $O(3)$ which preserves the
Euclidean distance $d\ell^2 = dx^2+dy^2+dz^2$.}

%%──────────────────────────────────────────────────────────────────
\subsection*{Key Results}

\thm{Invariance of the spacetime interval}{The spacetime interval between any two events is the same in all
inertial frames:
\begin{equation}
  \boxed{ds^2 = ds'^2}
  \label{eq:lec02:invariance}
\end{equation}
Equivalently, the proportionality factor $\lambda$ defined by
$ds^2 = \lambda(|\mathbf{v}|/c)\,ds'^2$ satisfies $\lambda = 1$.}

\thm{Lorentz group as the symmetry of the interval}{Lorentz transformations are the unique linear transformations
preserving $ds^2$.  Mathematically, they are defined as the group
$O(3,1)$, the set of matrices that preserve the indefinite quadratic
form $\eta_{\mu\nu}$ with signature $(+,-,-,-)$.  This is the
relativistic analogue of $O(3)$ preserving $\delta_{ij}$.}

%%──────────────────────────────────────────────────────────────────
\subsection*{Derivations}

%% DERIVATION 1: General argument for interval invariance
\subsubsection*{Proof (physical argument) that $ds^2 = ds'^2$ for
arbitrary events}

\paragraph{Setup.}
Let $S$ and $S'$ be two inertial frames, with $S'$ moving at velocity
$\mathbf{v}'$ relative to $S$.  Consider two events that are
infinitesimally close; denote their separation in $S$ by $ds$ and in
$S'$ by $ds'$.  The goal is to determine the relationship between
$ds^2$ and $ds'^2$ for \emph{arbitrary} events, not merely for
light-propagation events.

\paragraph{Step 1: most general ansatz.}
We write
\begin{equation}
  ds^2 = \lambda\bigl(t_1, \mathbf{r}_1,\,\mathbf{v}',\,
         \widehat{\mathbf{v}}'\bigr)\;ds'^2,
  \label{eq:lec02:ansatz0}
\end{equation}
where the proportionality factor $\lambda$ is, a priori, a function
of the location of the event $(t_1,\mathbf{r}_1)$, the magnitude of
the relative velocity $|\mathbf{v}'|$, and its direction
$\widehat{\mathbf{v}}'$.  (The infinitesimal separation $ds'^2$
itself already encodes the displacement $dt,\,d\mathbf{r}$, so no
additional dependence on $dt$ or $d\mathbf{r}$ is needed beyond
what is already in $ds'^2$.)

\paragraph{Step 2: eliminate dependence on location (homogeneity).}
Homogeneity of spacetime means no spacetime point is privileged:
the laws of physics — and hence the transformation factor — cannot
depend on where or when the events occur.  Therefore
\begin{equation}
  \lambda \;\text{ is independent of }\; t_1,\,\mathbf{r}_1.
  \label{eq:lec02:drop_r1}
\end{equation}

\paragraph{Step 3: eliminate dependence on the direction of $\mathbf{v}'$
(isotropy).}
Isotropy of space means no spatial direction is preferred.  The
magnitude of $ds^2$ (a scalar) cannot depend on which direction the
frame $S'$ moves relative to $S$.  Therefore
\begin{equation}
  \lambda = \lambda\!\left(\frac{|\mathbf{v}'|}{c}\right).
  \label{eq:lec02:drop_dir}
\end{equation}
The factor $1/c$ appears because $\lambda$ must be dimensionless and
$c$ is the only invariant speed in the theory.

After steps 2 and 3 the ansatz reduces to
\begin{equation}
  ds^2 = \lambda\!\left(\frac{|\mathbf{v}'|}{c}\right)\;ds'^2.
  \label{eq:lec02:ansatz}
\end{equation}

\paragraph{Step 4: three-frame consistency (key argument).}
Introduce a third inertial frame $S''$ moving at velocity
$\mathbf{v}''$ relative to $S$.  Applying \eqref{eq:lec02:ansatz}
twice:
\begin{align}
  ds^2   &= \lambda(|\mathbf{v}'|/c)\;ds'^2,
             \label{eq:lec02:eq1}\\
  ds^2   &= \lambda(|\mathbf{v}''|/c)\;ds''^2,
             \label{eq:lec02:eq2}\\
  ds'^2  &= \lambda(|\mathbf{v}|/c)\;ds''^2,
             \label{eq:lec02:eq3}
\end{align}
where $|\mathbf{v}|$ is the speed of $S''$ relative to $S'$.
Dividing \eqref{eq:lec02:eq1} by \eqref{eq:lec02:eq2}:
\begin{equation}
  \frac{ds'^2}{ds''^2}
  = \frac{\lambda(|\mathbf{v}''|/c)}{\lambda(|\mathbf{v}'|/c)}.
  \label{eq:lec02:div}
\end{equation}
But from \eqref{eq:lec02:eq3},
\begin{equation}
  \frac{ds'^2}{ds''^2} = \lambda(|\mathbf{v}|/c).
  \label{eq:lec02:from3}
\end{equation}
Equating \eqref{eq:lec02:div} and \eqref{eq:lec02:from3}:
\begin{equation}
  \lambda\!\left(\frac{|\mathbf{v}|}{c}\right)
  = \frac{\lambda(|\mathbf{v}''|/c)}{\lambda(|\mathbf{v}'|/c)}.
  \label{eq:lec02:functional}
\end{equation}

\paragraph{Step 5: the right-hand side cannot depend on the angle
between $\mathbf{v}'$ and $\mathbf{v}''$.}
By isotropy, $\lambda$ depends only on the \emph{magnitude} of its
argument, never on a direction.  The right-hand side of
\eqref{eq:lec02:functional} involves only the magnitudes
$|\mathbf{v}'|$ and $|\mathbf{v}''|$, so it carries no angular
information.  The left-hand side, however, contains
$|\mathbf{v}|$ — the speed of $S''$ relative to $S'$ — which in
general \emph{does} depend on the angle $\theta$ between $\mathbf{v}'$
and $\mathbf{v}''$.  (Concretely: if $\mathbf{v}' \parallel \mathbf{v}''$
then $|\mathbf{v}| = 0$; if $\mathbf{v}' \perp \mathbf{v}''$ then
$|\mathbf{v}| \neq 0$.)

Equation \eqref{eq:lec02:functional} therefore requires the left-hand
side to be independent of $\theta$, which is possible only if
$\lambda$ is \emph{constant} (independent of its argument):
\begin{equation}
  \lambda = \text{const.}
  \label{eq:lec02:const}
\end{equation}

\paragraph{Step 6: the constant must equal 1.}
Substitute $\lambda = \text{const}$ back into
\eqref{eq:lec02:functional}:
\begin{equation}
  \text{const} = \frac{\text{const}}{\text{const}}
  \;\implies\;
  \text{const}^2 = \text{const}
  \;\implies\;
  \text{const} = 1.
  \label{eq:lec02:const1}
\end{equation}
(The value $\text{const} = 0$ is excluded because $ds^2 = 0$ for all
events would be unphysical.)

\paragraph{Conclusion.}
\begin{equation}
  ds^2 = ds'^2
  \qquad \text{for any two infinitesimally separated events.}
  \label{eq:lec02:infinitesimal-invariance}
\end{equation}
By integration (summing infinitesimal intervals), the result
extends to any finite pair of events:
\begin{equation}
  \Delta s^2 = \Delta s'^2.
  \label{eq:lec02:finite-invariance}
\end{equation}

\paragraph{Remark on rigor.}
Smolkin emphasized explicitly that this is a \emph{physical argument},
not a rigorous mathematical proof.  The extension from the
infinitesimal to the finite case rests on the fact that we showed the
proportionality factor equals~1 \emph{at every spacetime point}, so
integrating up the infinitesimal intervals cannot introduce any
discrepancy.  Formally constructing the Lorentz group from the
requirement of interval preservation (the mathematician's route) makes
this rigorous but takes the interval invariance as an axiom rather
than deriving it.

%% DERIVATION 2: Light signals have ds^2 = 0 (recap from Lec01)
\subsubsection*{Recap: light signals satisfy $ds^2 = 0$ in every frame}

For a signal propagating at the maximum speed $c$ in frame $S$,
$|d\mathbf{r}| = c\,|dt|$, hence
\begin{equation}
  ds^2 = c^2\,dt^2 - d\mathbf{r}^2 = 0.
  \label{eq:lec02:null-interval}
\end{equation}
The principle of relativity demands that the same signal satisfies
$|d\mathbf{r}'| = c\,|dt'|$ in every frame $S'$ (the maximum speed
is a law of nature, hence frame-independent), so $ds'^2 = 0$.  This
is the special case $ds^2 = ds'^2 = 0$ that motivated the general
invariance result.

%% DERIVATION 3: Analogy with Euclidean rotations
\subsubsection*{Analogy: orthogonal group $O(3)$ preserves Euclidean
distance}

The Euclidean distance
$d\ell^2 = dx^2 + dy^2 + dz^2$
is preserved by all \emph{orthogonal} transformations $R \in O(3)$:
\begin{equation}
  d\ell^2 = (R\,d\mathbf{r})^T(R\,d\mathbf{r})
           = d\mathbf{r}^T R^T R\,d\mathbf{r}
           = d\mathbf{r}^T d\mathbf{r}
           = d\ell^2
  \quad (R^T R = \mathbf{1}).
  \label{eq:lec02:orthogonal-invariance}
\end{equation}
Lorentz transformations $\Lambda$ play the identical role for the
indefinite form
$ds^2 = c^2 dt^2 - dx^2 - dy^2 - dz^2 = \eta_{\mu\nu}\,dx^\mu\,dx^\nu$:
\begin{equation}
  \Lambda^T \eta\,\Lambda = \eta,
  \qquad
  \eta = \mathrm{diag}(+1,-1,-1,-1).
  \label{eq:lec02:lorentz-metric-condition}
\end{equation}
The group of such $\Lambda$ is $O(3,1)$, the Lorentz group.

%%──────────────────────────────────────────────────────────────────
\subsection*{Examples}

\ex{Why the angle between $\mathbf{v}'$ and $\mathbf{v}''$
matters for $|\mathbf{v}|$}{Let $|\mathbf{v}'| = |\mathbf{v}''| = V$.
\begin{itemize}
  \item If $\mathbf{v}'' \parallel \mathbf{v}'$ (same direction):
    $S''$ moves at speed $V$ in the same direction as $S'$; both
    frames co-move, so $|\mathbf{v}| = 0$.
  \item If $\mathbf{v}'' \perp \mathbf{v}'$:
    $S''$ moves orthogonally to $S'$; the relative speed
    $|\mathbf{v}| \neq 0$.
\end{itemize}
Thus the left-hand side of \eqref{eq:lec02:functional} varies with
the angle between $\mathbf{v}'$ and $\mathbf{v}''$, while the
right-hand side does not — contradicting \eqref{eq:lec02:functional}
unless $\lambda$ is constant.}

\ex{Integers under addition and the Lorentz group}{Just as $(\mathbb{Z},+)$ is a group with identity $0$ and inverse
$-n$ for each $n$, the Lorentz group is a set of transformations
closed under composition, with the identity transformation (no boost,
no rotation) and with an inverse for each element.  The specific
feature distinguishing the Lorentz group from $O(3)$ is the
indefinite signature of the preserved quadratic form.}

%%──────────────────────────────────────────────────────────────────
\subsection*{Connections}

\begin{itemize}
  \item \textbf{Lecture 01.}  The starting point is the Lecture~01
    result that Maxwell's equations are inconsistent with the Galilean
    group.  The natural question — what group \emph{does} leave
    Maxwell's equations invariant? — motivates the construction of
    the Lorentz group from the invariance of $ds^2$.
  \item \textbf{Euclidean geometry / $O(3)$.}  Smolkin used the
    orthogonal group $O(3)$ as the canonical example of a group
    defined by preservation of a quadratic form ($d\ell^2$).  The
    Lorentz group $O(3,1)$ is its spacetime analogue.
  \item \textbf{Next lecture.}  Interval invariance will be taken as
    the axiom from which the explicit form of the Lorentz
    transformation is derived.  Smolkin stated: ``invariance of the
    interval is sufficient to uniquely determine the transformation.''
  \item \textbf{Landau \& Lifshitz, Vol.~2, §2.}  The argument
    presented here (homogeneity $+$ isotropy $+$ three-frame
    consistency $\Rightarrow$ $\lambda = 1$) follows the L\&L
    treatment closely, though Smolkin noted it should be read as
    physical motivation rather than a rigorous proof.
\end{itemize}

%%──────────────────────────────────────────────────────────────────
\subsection*{To Remember}

\begin{itemize}
  \item The spacetime interval is
    $ds^2 = c^2\,dt^2 - d\mathbf{r}^2$; it is a Lorentz
    \emph{scalar} — the same in every inertial frame.
  \item The argument for $\lambda = 1$ uses three ingredients:
    (1)~homogeneity eliminates dependence on the location of the
    events; (2)~isotropy eliminates dependence on the direction of
    $\mathbf{v}'$, leaving $\lambda = \lambda(|\mathbf{v}'|/c)$;
    (3)~three-frame consistency forces $\lambda$ to be constant,
    and the constant must be~1.
  \item The speed $c$ enters as the unique invariant speed scale that
    makes $\lambda$ dimensionless; its invariance is a consequence of
    the principle of relativity applied to the maximum-speed signal,
    not an independent postulate at this stage.
  \item Lorentz transformations are the group $O(3,1)$, defined by
    $\Lambda^T \eta \Lambda = \eta$ with
    $\eta = \mathrm{diag}(+1,-1,-1,-1)$.  They play the same role
    for $ds^2$ as orthogonal matrices play for $d\ell^2$.
  \item The argument presented is physical motivation, not a
    mathematically rigorous proof.  The mathematician's route takes
    interval invariance as an axiom and defines the Lorentz group
    directly as the symmetry group of the quadratic form.
\end{itemize}

%%──────────────────────────────────────────────────────────────────
